% !TeX root = ./../main.tex

\chapter{Uczenie głębokie jako narzędzie redukcji błędów}\label{chapter:dl_error_reduction}

\agentinfo{
Rozdział opisuje użycie generatywnej sieci przeciwstawnej pix2pix do redukcji błędu w obrazach kodowanych metodą LPIQE, samodzielnie oraz w połączeniu z PDU.
Praca nie została opublikowana w czasopiśmie; istnieje jako preprint z 2023 r. (Preprints.org) i tak należy ją cytować.
Docelowa objętość rozdziału: około 10 stron.
Źródła:
manuskrypt w szablonie czasopisma Quantum (wersja po korekcie językowej, traktować jako nadrzędną) \agentref{bib:werner_deep_learning_manuscript} \agentref{artykuły/Quantum_journal_GAN__Copy_/article.tex},
ten sam manuskrypt w szablonie MDPI \agentref{artykuły/Energies_article/template.tex},
rysunki \agentref{artykuły/Quantum_journal_GAN__Copy_/imgs}.
Podstawy sieci GAN i pix2pix są w zapleczu teoretycznym \agentref{chapters/theoretical_background_06_ML.tex}.
}

\section{Cel i zakres rozdziału}\label{sec:dl:scope}
\agentinfo{
Pytanie badawcze rozdziału: PB4.
Odniesienie do roboczej tezy T3.
}

\section{Redukcja błędu jako translacja obraz--obraz}\label{sec:dl:concept}
\agentinfo{
Sformułowanie zadania: wejściem jest obraz zrekonstruowany z pomiarów (surowy albo po PDU), wyjściem obraz zbliżony do oryginału.
Uzasadnienie: zniekształcenie jest systematyczne i zależne od urządzenia, więc może być wyuczone.
Miejsce sieci w potoku: po rekonstrukcji i po PDU, w całości po stronie klasycznej.
}

\section{Architektura sieci}\label{sec:dl:architecture}
\agentinfo{
Generator U-Net (liczba parametrów, bloki kodujące i dekodujące, połączenia skrótowe), dyskryminator PatchGAN, funkcja straty z wagami 1 i 100, optymalizator Adam z parametrami.
Skalowanie obrazów do 256x256 i z powrotem; omówić konsekwencje.
Schematy architektury wykonać na nowo w TikZ zamiast przenosić rastrowe rysunki z manuskryptu.
}

\section{Protokół eksperymentalny}\label{sec:dl:protocol}
\agentinfo{
Cztery eksperymenty:
E1: obrazy 16x16 z symulatora, odniesieniem obrazy po PDU;
E2: obrazy 8x8 z rzeczywistego urządzenia, odniesieniem obrazy po PDU;
E3: obrazy 16x16 z symulatora po PDU, odniesieniem obrazy oryginalne;
E4: obrazy 8x8 z rzeczywistego urządzenia po PDU, odniesieniem obrazy oryginalne.
Zbiory: 631 obrazów (symulator) i 61 obrazów (urządzenie IBMQ Nairobi); liczba epok 250 i 1000.
Opisać podział na zbiór uczący i testowy oraz kroki odtworzenia eksperymentu.
W manuskrypcie tytuły podrozdziałów wyników i lista eksperymentów nie są w pełni spójne; uzgodnić numerację z danymi przed opisem.
}

\section{Wyniki}\label{sec:dl:results}

\subsection{Obrazy z symulatora}\label{subsec:dl:results:sim}
\agentinfo{Eksperymenty E1 i E3.}

\subsection{Obrazy z rzeczywistego urządzenia}\label{subsec:dl:results:real}
\agentinfo{
Eksperymenty E2 i E4.
Wynik kluczowy do zweryfikowania w danych: MSE rzędu 0,00014 i odchylenie standardowe błędu 0,00993 dla E4.
}

\subsection{Porównanie: wynik surowy, PDU, sieć, PDU z siecią}\label{subsec:dl:results:comparison}
\agentinfo{Tabela zbiorcza miar dla wszystkich wariantów.}

\section{Dyskusja}\label{sec:dl:discussion}
\agentinfo{
Odpowiedź na PB4.
Zagrożenia trafności, które trzeba omówić wprost:
bardzo mały zbiór z rzeczywistego urządzenia (61 obrazów) i ryzyko przeuczenia;
ryzyko, że sieć odtwarza rozkład danych uczących zamiast usuwać szum (obraz wiarygodny wizualnie, lecz niezgodny z pomiarem);
zależność wytrenowanej sieci od urządzenia i od kalibracji;
koszt treningu sieci o ponad 54 milionach parametrów dla obrazów 8x8 i 16x16;
skalowanie na większe obrazy.
Jeśli dane na to pozwalają, rozważyć eksperymenty uzupełniające: walidację krzyżową, test na obrazach spoza rozkładu, prostsze modele odniesienia (do decyzji autora).
}

\section{Podsumowanie}\label{sec:dl:summary}
